\documentclass[12pt]{article}
\usepackage[utf8]{inputenc}
\usepackage{amsmath, amssymb, amsthm}
\usepackage{geometry}

% Set margins to match a standard homework document
\geometry{a4paper, margin=1in}

\title{Challenge Problem Set 2}
\author{Math 292 Spring 2026 \\ Rutgers University}
\date{}

\begin{document}

\maketitle

\noindent \textbf{Exercise 4.} Let us fix the initial data to be
\[
    x_{0}=1, \quad u_{0}=0, \quad y_{0}=1, \quad \text{and} \quad v_{0}=\sqrt{3}.
\]
Let $\mathbf{x}(t)$ be the unique solution of $x''(t)=-x(t)$ with $x(0)=(x_{0},y_{0})$ and $x'(0)=(u_{0},v_{0})$ for these values of the initial data.

\begin{enumerate}

    \item[(a)] Find $R_{2}=\max_{t\in\mathbb{R}}\|\mathbf{x}(t)\|$, and find all $t$ for which $\|\mathbf{x}(t)\|=R_{2}$.

          \textbf{Solution:}

          1. From the previous parts, we found that that the $||\vec{x}(t)||$ is maximized when
          \[
              \vec{x}(t) = \frac{1}{\sqrt{\lambda{}_2}}\vec{u}_2
          \]

          2. Thus, we must solve for the eigenvalues and eigenvectors of the matrix M with the initial data given above.
          \begin{align*}
              M & = (A^{-1})^T(A^{-1})                                         \\
                & = \frac{1}{3}
              \begin{bmatrix} \sqrt{3} & -1 \\ 0 & 1 \end{bmatrix}
              \begin{bmatrix} \sqrt{3} & 0 \\ -1 & 1 \end{bmatrix}             \\
                & = \frac{1}{3} \begin{bmatrix} 4 & -1 \\ -1 & 1 \end{bmatrix}
          \end{align*}

          3. To do this, we first solve the characteristic polynomial, letting $\gamma{} = 3\lambda{}$
          \begin{align*}
              det(A - \lambda{}I) & = \frac{1}{3}(4-\gamma{})(1-\gamma{}) - (-1)(-1) \\
                                  & = \gamma{}^2 - 5\gamma{} + 3                     \\
                                  & = 0
          \end{align*}

          4. We can now solve for the eigenvalues using the quadratic formula
          \begin{align*}
              \gamma{} & = \frac{5 \pm{} \sqrt{(-5)^2 - 4(1)(3)}}{2(1)}      \\
              \gamma{} & = \frac{5 \pm{} \sqrt{13}}{2}                       \\
                       & \implies{} \lambda{}  = \frac{5 \pm{} \sqrt{13}}{6}
          \end{align*}

          5. Remember that $\lambda{}_1 > \lambda{}_2$
          \begin{align*}
              \lambda{}_1 & = \frac{5 + \sqrt{13}}{6} \\
              \lambda{}_2 & = \frac{5 - \sqrt{13}}{6}
          \end{align*}

          6. With this, we can solve for $R_2$, take the norm of both sides of the equation from step one
          \begin{align*}
              ||\vec{x}(t)|| & = \frac{1}{\sqrt{\lambda{}_2}}||\vec{u}_2||   \\
              R_2            & = \sqrt{\frac{6}{5-\sqrt{13}}}~\approx{2.074}
          \end{align*}

          7. To solve for the time, we need the eigenvectors of M
          \begin{align*}
              (M-\lambda_{2}I)\vec{u} & = \begin{bmatrix} \frac{4}{3}-\frac{5 - \sqrt{13}}{6} & -\frac{1}{3} \\ -\frac{1}{3} & \frac{1}{3} - \frac{5 - \sqrt{13}}{6} \end{bmatrix} \vec{u} = 0
          \end{align*}

          8. Since we know this matrix above must be singular, we can solve for the eigenvector using the first row of the matrix. Let x, y be the coordinates of vector $\vec{u}_1$
          \begin{align*}
              (\frac{4}{3} - \frac{5-\sqrt{13}}{6})x -\frac{1}{3}y & = 0            \\
              (\frac{3+\sqrt{13}}{6})x                             & = \frac{2}{6}y \\
              (3+\sqrt{13})x                                       & = 2y
          \end{align*}

          9. Thus if we let $x = y = 1$, we get the eigenvector
          \[
              \vec{u}_2 = \begin{bmatrix}
                  2 \\ 3 + \sqrt{13}
              \end{bmatrix}
          \]

          10. From an earlier part, recall that $A <cos(t), sin(t)> = \pm{\frac{1}{\sqrt{\lambda{}_2}}}\vec{u_2}$, expanding we get
          \begin{align*}
              \begin{bmatrix} cos(t) \\ cos(t) + \sqrt{3}sin(t) \end{bmatrix} = \frac{1}{\sqrt{\lambda_{2}}}\begin{bmatrix} 2 \\ 3 + \sqrt{13} \end{bmatrix}
          \end{align*}

          11. We can cancel out the scalar by dividing across row 2 by row 1
          \begin{align*}
              \frac{3+\sqrt{13}}{2}         & = \frac{cos(t) + \sqrt{3}sin(t)}{cos(t)}               \\
              \frac{3+\sqrt{13}}{2}         & = 1 + \sqrt{3}tan(t)                                   \\
              \frac{1+\sqrt{13}}{2\sqrt{3}} & = tan(t)                                               \\
              t                             & = arctan(\frac{1+\sqrt{13}}{2\sqrt{3}}) \approx{0.926}
          \end{align*}

    \item[(b)] Find $R_{1}=\min_{t\in\mathbb{R}}\|\mathbf{x}(t)\|$, and find all $t$ for which $\|\mathbf{x}(t)\|=R_{1}$.

          \textbf{Solution:}

          1. Using the minimizer formula we found from previous exercise, we have
          \begin{align*}
              \vec{x}(t) = \frac{1}{\sqrt{\lambda{}_1}}\vec{u}_1
          \end{align*}

          2. Take the norm of both sides to get the minimum value
          \begin{align*}
              ||\vec{x}(t)|| & = \frac{1}{\sqrt{\lambda{}_1}}||\vec{u}_1||     \\
              R_1            & = \sqrt{\frac{6}{3+\sqrt{13}}} ~ \approx{0.835}
          \end{align*}

          3. Eigenvectors are solved in the same way as we did in part a
          \begin{align*}
              ( M-\lambda_{1}I)\vec{u}_1 = \begin{bmatrix} \frac{8}{6} - \frac{5+\sqrt{13}}{6} & \frac{-1}{3} \\ \frac{-1}{3} & \frac{2}{6}-\frac{5+\sqrt{13}}{6} \end{bmatrix} \vec{u}_1 = 0
          \end{align*}

          4. Find the ratio between the coordinates again
          \begin{align}
              (\frac{8-5-\sqrt{13}}{6})x & = \frac{2}{6}y \\
              (3-\sqrt{13})x = 2y
          \end{align}

          5. Setting $x=y=1$, we have
          \[
              \vec{u}_1 = \begin{bmatrix} 2 \\ 3-\sqrt{13} \end{bmatrix}
          \]

          6. Expanding $A <cos(t), sin(t)> = \pm{\frac{1}{\sqrt{\lambda{}_2}}}\vec{u_2}$ we have
          \begin{align*}
              \begin{bmatrix} cos(t) \\ cos(t) + \sqrt{3}sin(t) \end{bmatrix} = \frac{1}{\sqrt{\lambda_{1}}}\begin{bmatrix} 2 \\ 3 - \sqrt{13} \end{bmatrix}
          \end{align*}

          7. Dividing row 2 by row 1, we get
          \begin{align*}
              \frac{3-\sqrt{13}}{2}         & = \frac{cos(t) + \sqrt{3}sin(t)}{cos(t)}                \\
              \frac{1-\sqrt{13}}{2\sqrt{3}} & = tan(t)                                                \\
              t                             & = arctan(\frac{1-\sqrt{13}}{2\sqrt{3}}) \approx{-0.645}
          \end{align*}

    \item[(c)] Sketch a plot of the ellipse traced out by $\mathbf{x}(t)$ as $t$ ranges over $\mathbb{R}$.

          \textbf{Solution:} Attached on last page.
\end{enumerate}


In solving the concrete problem in Exercise 4, we have encountered some important constructions in Linear Algebra. Introducing the numbers $\sigma_{1}$ and $\sigma_{2}$ through
\[
    \sigma_{1}=\frac{1}{\sqrt{\lambda_{1}}} \quad \text{and} \quad \sigma_{2}=\frac{1}{\sqrt{\lambda_{2}}}
\]
we can rewrite (0.9) as
\begin{equation}
    Av_{1}=\sigma_{1}u_{1} \quad \text{and} \quad Av_{2}=\sigma_{2}u_{2} \tag{0.10}
\end{equation}
where $\{u_{1},u_{2}\}$ and $\{v_{1},v_{2}\}$ are both orthonormal bases of $\mathbb{R}^{2}$. Since any vector $x\in\mathbb{R}^{2}$ can be written as
\[
    x=(x\cdot v_{1})v_{1}+(x\cdot v_{2})v_{2},
\]
applying $A$ to both sides, and using (0.10), we obtain
\begin{equation}
    Ax=\sigma_{1}(x\cdot v_{1})u_{1}+\sigma_{2}(x\cdot v_{2})u_{2}. \tag{0.11}
\end{equation}
It turns out that this gives a factorization of the general invertible $2\times2$ matrix into a product of simple matrices. This is the singular value decomposition.

\bigskip
\noindent \textbf{Exercise 5.} Introduce the matrices $V=[v_{1},v_{2}]$,
\[
    S=\begin{bmatrix}\sigma_{1}&0\\ 0&\sigma_{2}\end{bmatrix}
\]
and $U=[u_{1},u_{2}]$. (That is, the columns of $V$ are $v_{1}$ and $v_{2}$, and the columns of $U$ are $u_{1}$ and $u_{2}$.) Show that (0.11) is equivalent to
\[
    Ax=USV^{T}x
\]
for all $x$, and thus that the matrix $A$ can be factored as the product of the three matrices $U$, $S$ and $V$:
\[
    A=USV^{T}
\]

\textbf{Solution:}

1. It will be easier to start with matrix form and expand out one by one to get (0.11) to show equivalence:
\begin{align*}
    A\vec{x} & = USV^T\vec{x}                                                                                                                                                \\
             & = \begin{bmatrix} \vec{u}_1 & \vec{u}_2 \end{bmatrix} \begin{bmatrix}       \sigma_{1} & 0 \\ 0 & \sigma_{2}       \end{bmatrix}V^T\vec{x}                    \\
             & = \begin{bmatrix} \sigma_{1}\vec{u}_1 & \sigma_{2}\vec{u}_2\end{bmatrix} \begin{bmatrix} \vec{v}_1 \\ \vec{v}_2 \end{bmatrix}\vec{x}                          \\
             & = \begin{bmatrix} \sigma_{1}\vec{u}_1 & \sigma_{2}\vec{u}_2\end{bmatrix} \begin{bmatrix} \vec{x} \cdot{} \vec{v}_1 \\ \vec{x} \cdot{} \vec{v}_2 \end{bmatrix} \\
             & = \sigma_{1}(x\cdot v_{1})u_{1}+\sigma_{2}(x\cdot v_{2})u_{2} \tag{0.11}
\end{align*}

\hrulefill{}

Recall that a square matrix with orthonormal columns is an orthogonal matrix, and that the inverse of an orthogonal matrix is its transpose. Hence $U$ and $V$ are orthogonal matrices, and their inverses are their transposes. The only $2\times2$ orthogonal matrices are rotations and reflections, so $U$ and $V$ both represent simple geometric transformations. The matrix $S$ is even simpler: It represents a scaling transformation: It simply rescales the $x$ and $y$ coordinates by the factors $\sigma_{1}$ and $\sigma_{2}$ respectively.

Thus, every invertible linear transformation of $\mathbb{R}^{2}$ into itself is the product of an orthogonal transformation, followed by a scaling transformation, followed by another orthogonal transformation.

Here is one consequence of this important fact. Orthogonal transformations (rotations and reflections) do not affect area. Hence the effect on area of the linear transformation given by $A$ comes entirely from the scale transformation in the middle. Thus, if $R$ is a region in the plane, and $A(R)$ is its image under $A$, the area of $A(R)$ is $\sigma_{1}\sigma_{2}$ times the area of $R$, and it is easy to see that $\sigma_{1}\sigma_{2}=|\det(A)|$.

All of this generalizes to matrices of arbitrary size, but in this workshop we will stay with 2 dimensions, and thoroughly deal with this case.

The final problem is to compute a singular value decomposition. Now that we know that such a decomposition exists, we can efficiently find it as follows. Starting from $A=USV^{T}$, we write
\[
    A^{T}A=(USV^{T})^{T}USV^{T}=VSU^{T}USV^{T}=VS^{2}V^{T}
\]
where we have used the fact that for any matrix product $XY$, $(XY)^{T}=Y^{T}X^{T}$ and the fact that since $U^{T}$ is the inverse of $U$, since $U$ is orthogonal. Then since $V^{T}$ is the inverse of $V$, again, since $V$ is orthogonal,
\[
    V^{T}(A^{T}A)V=V^{-1}(A^{T}A)V=S^{2}.
\]
Thus, $V$ diagonalizes the matrix $A^{T}A$. This means the columns of $V$ are an orthonormal basis of eigenvectors of $A^{T}A$, and the corresponding eigenvalues are the squares of the diagonal entries of $S$. So to find $\{v_{1},v_{2}\}$ and $\{\sigma_{1},\sigma_{2}\}$, diagonalize $A^{T}A$. Finally, to find $\{u_{1},u_{2}\}$, use (0.10).

\bigskip
\noindent \textbf{Exercise 6.} Let $A=\begin{bmatrix}11&-5\\ 2&-10\end{bmatrix}.$
Compute a singular value decomposition of $A$. That is, find orthogonal matrices $U$ and $V$, and a diagonal matrix $S$ with positive entries (these are the singular values) so that
\[
    A=USV^{T}.
\]
All of the terms in the singular value decomposition have a direct meaning for the solution of the differential equation $x''=-x$ when the matrix $A$ is given by the initial data, as explained above. This is one way to think about where the singular value decomposition "comes from".

\textbf{Solution:}

1. To find the SVD of A, we begin by calculate $A^TA$
\begin{align*}
    A^TA & = \begin{bmatrix} 11 & 2 \\ -5 & -10 \end{bmatrix}  \begin{bmatrix} 11 & -5 \\ 2 & -10 \end{bmatrix} \\
         & = \begin{bmatrix} 125 & -75 \\ -75 & 125 \end{bmatrix}
\end{align*}

2. Then we solve the characteristic polynomial from $det(A^TA-\lambda{}I) = 0$ for it's eigenvalues
\begin{align*}
    (125-\lambda{})(125-\lambda{})-(-75*-75) & = \lambda{}^2 - 250\lambda{} + 10000 = 0  \\
                                             & = (\lambda{} - 50)(\lambda{} - 200) = 0   \\ \\
    \lambda_{1} = 200                        & ~,~                     \lambda_{2}  = 50
\end{align*}

3. Thus, we have the our two singular values and matrix S
\begin{align*}
    \sigma_{1} & = \sqrt{200} = 10\sqrt{2}        \\
    \sigma_{2} & = \sqrt{50} = 5\sqrt{2}          \\
    S          & = \begin{bmatrix}
                       10\sqrt{2} & 0 \\ 0 & 5\sqrt{2}
                   \end{bmatrix}
\end{align*}

4. Now we solve for the eigenvectors of $A^TA$ which is equal to V
\begin{itemize}f
    \item For $\lambda_{1} = 200$:
          \begin{align*}
              \begin{bmatrix} -75 & -75 \\ -75 & -75 \end{bmatrix}\begin{bmatrix} x \\ y \end{bmatrix} = 0 \\
              -75x - 75y = 0 \implies{x = -y}
          \end{align*}

          Normalizing, we have our first eigenvector
          \[
              \vec{v}_1 = \begin{bmatrix}\frac{1}{\sqrt{2}} \\ -\frac{1}{\sqrt{2}} \end{bmatrix}
          \]
    \item For $\lambda_{2} = 50$:
          \begin{align*}
              \begin{bmatrix} 75 & -75 \\ -75 & 75 \end{bmatrix}\begin{bmatrix} x \\ y \end{bmatrix} = 0 \\
              75x - 75y = 0 \implies{x = y}
          \end{align*}

          Normalizing again, we now have our second eigenvector
          \[
              \vec{v}_2 = \begin{bmatrix}\frac{1}{\sqrt{2}} \\ \frac{1}{\sqrt{2}} \end{bmatrix}
          \]
\end{itemize}

5. Thus our matrix V is
\begin{equation*}
    V = \begin{bmatrix}\frac{1}{\sqrt{2}} & \frac{1}{\sqrt{2}} \\ -\frac{1}{\sqrt{2}} & \frac{1}{\sqrt{2}} \end{bmatrix}
\end{equation*}

6. To get the matrix U, we rewrite the equation (0.10)
\begin{itemize}
    \item For $\vec{u}_1$:
          \begin{align*}
              \vec{u}_1 & = \frac{1}{\sigma_{1}} A \vec{v}_1                                                                                                                                            \\
                        & = \frac{1}{10\sqrt{2}} \begin{bmatrix} 11 & -5 \\ 2 & -10 \end{bmatrix} \begin{bmatrix} \frac{1}{\sqrt{2}} \\ -\frac{1}{\sqrt{2}} \end{bmatrix} \\
                        & = \frac{1}{20} \begin{bmatrix} 16 \\ 12 \end{bmatrix}                                                                                                                         \\
                        & = \begin{bmatrix} \frac{4}{5} \\ \frac{3}{5} \end{bmatrix}
          \end{align*}
    \item For $\vec{u}_2$:
          \begin{align*}
              \vec{u}_2 & = \frac{1}{\sigma_{2}} A \vec{v}_2                                                                                                                                          \\
                        & = \frac{1}{5\sqrt{2}} \begin{bmatrix} 11 & -5 \\ 2 & -10 \end{bmatrix} \begin{bmatrix} \frac{1}{\sqrt{2}} \\ \frac{1}{\sqrt{2}} \end{bmatrix} \\
                        & = \frac{1}{10} \begin{bmatrix} 6 \\ 8 \end{bmatrix}                                                                                                                         \\
                        & = \begin{bmatrix} \frac{3}{5} \\ -\frac{4}{5} \end{bmatrix}
          \end{align*}
\end{itemize}

7. Finally, we can assemble the SVD
\begin{align}
    A & = USV^T                                                                                 \\
      & = \begin{bmatrix} \frac{4}{5} & \frac{3}{5} \\ \frac{3}{5} & -\frac{4}{5} \end{bmatrix}
    \begin{bmatrix} 10\sqrt{2} & 0 \\ 0 & 5\sqrt{2} \end{bmatrix}
    \begin{bmatrix} \frac{1}{\sqrt{2}} & -\frac{1}{\sqrt{2}} \\ \frac{1}{\sqrt{2}} & \frac{1}{\sqrt{2}} \end{bmatrix}
\end{align}


\end{document}